\documentclass[10pt,a4paper]{amsart}
\usepackage{amsmath,amssymb,amsthm,graphicx}
% Retained minimal preamble for the Stacks Project chapter coherent.tex (Stacks chapter, pre-delegation retained fragment for stacks-coherent).
% Verbatim excerpt of the chapter preamble (sources/stacks-spaces-descent/preamble.tex): the reference/slogan/history
% environment block (native lines 7-12), Xy-pic setup (native lines 14-19) and the native theorem declarations with
% the mathematical macros (native lines 154-217). Copied mechanically; nothing added, removed or redefined. The
% wrapper class and the amsmath/amssymb/amsthm/graphicx base packages are supplied by the assembler, so the
% class/xr-hyper/hyperref/externaldocument part of the native preamble is omitted.
% For dealing with references we use the comment environment
\usepackage{verbatim}
\newenvironment{reference}{\comment}{\endcomment}
%\newenvironment{reference}{}{}
\newenvironment{slogan}{\comment}{\endcomment}
\newenvironment{history}{\comment}{\endcomment}
% For commutative diagrams we use Xy-pic
\usepackage[all]{xy}

% We use 2cell for 2-commutative diagrams.
\xyoption{2cell}
\UseAllTwocells
% Theorem environments.
%
\theoremstyle{plain}
\newtheorem{theorem}[subsection]{Theorem}
\newtheorem{proposition}[subsection]{Proposition}
\newtheorem{lemma}[subsection]{Lemma}

\theoremstyle{definition}
\newtheorem{definition}[subsection]{Definition}
\newtheorem{example}[subsection]{Example}
\newtheorem{exercise}[subsection]{Exercise}
\newtheorem{situation}[subsection]{Situation}

\theoremstyle{remark}
\newtheorem{remark}[subsection]{Remark}
\newtheorem{remarks}[subsection]{Remarks}

\numberwithin{equation}{subsection}

% Macros
%
\def\lim{\mathop{\mathrm{lim}}\nolimits}
\def\colim{\mathop{\mathrm{colim}}\nolimits}
\def\Spec{\mathop{\mathrm{Spec}}}
\def\Hom{\mathop{\mathrm{Hom}}\nolimits}
\def\Ext{\mathop{\mathrm{Ext}}\nolimits}
\def\SheafHom{\mathop{\mathcal{H}\!\mathit{om}}\nolimits}
\def\SheafExt{\mathop{\mathcal{E}\!\mathit{xt}}\nolimits}
\def\Sch{\mathit{Sch}}
\def\Mor{\mathop{\mathrm{Mor}}\nolimits}
\def\Ob{\mathop{\mathrm{Ob}}\nolimits}
\def\Sh{\mathop{\mathit{Sh}}\nolimits}
\def\NL{\mathop{N\!L}\nolimits}
\def\CH{\mathop{\mathrm{CH}}\nolimits}
\def\proetale{{pro\text{-}\acute{e}tale}}
\def\etale{{\acute{e}tale}}
\def\QCoh{\mathit{QCoh}}
\def\Ker{\mathop{\mathrm{Ker}}}
\def\Im{\mathop{\mathrm{Im}}}
\def\Coker{\mathop{\mathrm{Coker}}}
\def\Coim{\mathop{\mathrm{Coim}}}

% Boxtimes
%
\DeclareMathSymbol{\boxtimes}{\mathbin}{AMSa}{"02}

%
% Macros for moduli stacks/spaces
%
\def\QCohstack{\mathcal{QC}\!\mathit{oh}}
\def\Cohstack{\mathcal{C}\!\mathit{oh}}
\def\Spacesstack{\mathcal{S}\!\mathit{paces}}
\def\Quotfunctor{\mathrm{Quot}}
\def\Hilbfunctor{\mathrm{Hilb}}
\def\Curvesstack{\mathcal{C}\!\mathit{urves}}
\def\Polarizedstack{\mathcal{P}\!\mathit{olarized}}
\def\Complexesstack{\mathcal{C}\!\mathit{omplexes}}
% \Pic is the operator that assigns to X its picard group, usage \Pic(X)
% \Picardstack_{X/B} denotes the Picard stack of X over B
% \Picardfunctor_{X/B} denotes the Picard functor of X over B
\def\Pic{\mathop{\mathrm{Pic}}\nolimits}
\def\Picardstack{\mathcal{P}\!\mathit{ic}}
\def\Picardfunctor{\mathrm{Pic}}
\def\Deformationcategory{\mathcal{D}\!\mathit{ef}}
\usepackage{cleveref}
\crefname{originalproposition}{original proposition}{original propositions}
\Crefname{originalproposition}{Original proposition}{Original propositions}
\makeatletter
\newlabel{constructions-lemma-proj-sheaves}{{01M7}{source}{}{sourcepointer-1}{}}
\newlabel{constructions-lemma-proj-sheaves@cref}{{[originalproposition][0][]01M7}{[1][0][]source}{}{}{}}
\makeatother
\crefname{originalproposition}{original proposition}{original propositions}
\Crefname{originalproposition}{Original proposition}{Original propositions}
\makeatletter
\newlabel{homology-lemma-kernel-exact-functor}{{02MQ}{source}{}{sourcepointer-2}{}}
\newlabel{homology-lemma-kernel-exact-functor@cref}{{[originalproposition][0][]02MQ}{[1][0][]source}{}{}{}}
\makeatother
\crefname{originalproposition}{original proposition}{original propositions}
\Crefname{originalproposition}{Original proposition}{Original propositions}
\makeatletter
\newlabel{homology-lemma-serre-subcategory-is-kernel}{{02MS}{source}{}{sourcepointer-3}{}}
\newlabel{homology-lemma-serre-subcategory-is-kernel@cref}{{[originalproposition][0][]02MS}{[1][0][]source}{}{}{}}
\makeatother
\crefname{originalproposition}{original proposition}{original propositions}
\Crefname{originalproposition}{Original proposition}{Original propositions}
\makeatletter
\newlabel{lemma-coherent-on-proj-general}{{0B5Q}{source}{}{sourcepointer-4}{}}
\newlabel{lemma-coherent-on-proj-general@cref}{{[originalproposition][0][]0B5Q}{[1][0][]source}{}{}{}}
\makeatother
\crefname{originalproposition}{original proposition}{original propositions}
\Crefname{originalproposition}{Original proposition}{Original propositions}
\makeatletter
\newlabel{lemma-recover-tail-graded-module-general}{{0B5R}{source}{}{sourcepointer-5}{}}
\newlabel{lemma-recover-tail-graded-module-general@cref}{{[originalproposition][0][]0B5R}{[1][0][]source}{}{}{}}
\makeatother
\crefname{originalproposition}{original proposition}{original propositions}
\Crefname{originalproposition}{Original proposition}{Original propositions}
\makeatletter
\newlabel{properties-lemma-proj-quasi-coherent}{{0AG5}{source}{}{sourcepointer-6}{}}
\newlabel{properties-lemma-proj-quasi-coherent@cref}{{[originalproposition][0][]0AG5}{[1][0][]source}{}{}{}}
\makeatother
\crefname{originalproposition}{original proposition}{original propositions}
\Crefname{originalproposition}{Original proposition}{Original propositions}
\makeatletter
\newlabel{proposition-coherent-modules-on-proj}{{0BXD}{source}{}{sourcepointer-7}{}}
\newlabel{proposition-coherent-modules-on-proj@cref}{{[originalproposition][0][]0BXD}{[1][0][]source}{}{}{}}
\makeatother
\title{Coherent modules on Proj of a general Noetherian graded ring: for A Noetherian graded and X = Proj(A), the functor M -> M\textasciicircum{}tilde induces an equivalence Mod\textasciicircum{}fg\_A / Mod\textasciicircum{}fg\_\{A,irrelevant\} -> Coh(O\_X) whose quasi-inverse is F |-> bigoplus\_\{n >= 0\} Gamma(X, F(n)).}
\author{The Stacks Project Authors}
\date{stacks-project:coherent.tex:local-snapshot-20261012; GNU FDL 1.2 or later}
\begin{document}
\maketitle
\noindent\textit{Editorial note.} Statement span is lines 4442-4453 of sources/stacks-coherent/pinned-original-source.tex; proof span is lines 4455-4527, the complete author proof (71 body lines).\par\medskip
\noindent\textit{Editorial note.} Navigation stubs: 7 cited-but-not-included labels are supplied as non-load-bearing pointers with their pinned tags where the snapshot tag index has them.\par\medskip
\noindent\textit{Editorial note.} extra\_packages is empty on purpose: the retained preamble already loads xy with options and a second declaration clashes in the assembler wrapper.\par\medskip
\noindent\textit{Editorial note.} Text is retained verbatim, including the author's own typos and placeholders.\par\medskip
\noindent\textit{Editorial note.} Full derivation in span. Depends on lemma-recover-tail-graded-module-general and lemma-coherent-on-proj-general (same chapter; the former is installed in this wave as L16356).\par\medskip
\noindent\textit{Editorial note.} Span computed from the \textbackslash{}label marker through the enclosing \textbackslash{}end\{proposition\}; precard\_gate RESULT=OK (probe L9999).\par\medskip
\noindent\textit{Allowed standard prerequisites (editorial).}\begin{itemize}
\item Explicit prerequisite boundary. See why-field; every pointer named here occurs verbatim in the retained statement/proof text.
\end{itemize}\medskip
\ifdefined\hypertarget\hypertarget{sourcepointer-1}{}\fi
\noindent\textit{Source reference.} Original proposition 01M7: sources/stacks-coherent/pinned-original-source.tex, source lines 4527--4527. Cited by the retained statement or proof but not part of the collected spans; supplied only as a navigation pointer. The argument does not depend on its content.\par\medskip
\ifdefined\hypertarget\hypertarget{sourcepointer-2}{}\fi
\noindent\textit{Source reference.} Original proposition 02MQ: sources/stacks-coherent/pinned-original-source.tex, source lines 4527--4527. Cited by the retained statement or proof but not part of the collected spans; supplied only as a navigation pointer. The argument does not depend on its content.\par\medskip
\ifdefined\hypertarget\hypertarget{sourcepointer-3}{}\fi
\noindent\textit{Source reference.} Original proposition 02MS: sources/stacks-coherent/pinned-original-source.tex, source lines 4527--4527. Cited by the retained statement or proof but not part of the collected spans; supplied only as a navigation pointer. The argument does not depend on its content.\par\medskip
\ifdefined\hypertarget\hypertarget{sourcepointer-4}{}\fi
\noindent\textit{Source reference.} Original proposition 0B5Q: sources/stacks-coherent/pinned-original-source.tex, source lines 4237--4255. Cited by the retained statement or proof but not part of the collected spans; supplied only as a navigation pointer. The argument does not depend on its content.\par\medskip
\ifdefined\hypertarget\hypertarget{sourcepointer-5}{}\fi
\noindent\textit{Source reference.} Original proposition 0B5R: sources/stacks-coherent/pinned-original-source.tex, source lines 4298--4313. Cited by the retained statement or proof but not part of the collected spans; supplied only as a navigation pointer. The argument does not depend on its content.\par\medskip
\ifdefined\hypertarget\hypertarget{sourcepointer-6}{}\fi
\noindent\textit{Source reference.} Original proposition 0AG5: sources/stacks-coherent/pinned-original-source.tex, source lines 4527--4527. Cited by the retained statement or proof but not part of the collected spans; supplied only as a navigation pointer. The argument does not depend on its content.\par\medskip
\ifdefined\hypertarget\hypertarget{sourcepointer-7}{}\fi
\noindent\textit{Source reference.} Original proposition 0BXD: sources/stacks-coherent/pinned-original-source.tex, source lines 4168--4180. Cited by the retained statement or proof but not part of the collected spans; supplied only as a navigation pointer. The argument does not depend on its content.\par\medskip
\begin{proposition}
\label{proposition-coherent-modules-on-proj-general}
Let $A$ be a Noetherian graded ring. Set $X = \text{Proj}(A)$. The functor
$M \mapsto \widetilde M$ induces an equivalence
$$
\text{Mod}^{fg}_A/\text{Mod}^{fg}_{A, irrelevant}
\longrightarrow
\textit{Coh}(\mathcal{O}_X)
$$
whose quasi-inverse is given by
$\mathcal{F} \longmapsto \bigoplus_{n \geq 0} \Gamma(X, \mathcal{F}(n))$.
\end{proposition}

\begin{proof}
We urge the reader to read the proof in the case where $A$
is generated in degree $1$ first, see
Proposition \ref{proposition-coherent-modules-on-proj}.
Let $f_1, \ldots, f_r \in A$ be homogeneous elements of positive degree
which generate $A$ over $A_0$. Let $d$ be the lcm of the degrees
$d_i$ of $f_i$. Let $M$ be a finite $A$-module.
Let us show that $\widetilde{M}$ is zero if and
only if $M$ is an irrelevant graded $A$-module (as defined above
the statement of the proposition). Namely, let $x \in M$ be a
homogeneous element.
Choose $k \in \mathbf{Z}$ sufficiently small and let
$N \to N'$ and $M \to N'$ be as in
Lemma \ref{lemma-recover-tail-graded-module-general}.
We may also pick $l$ sufficiently large such that
$M_n \to N_n$ is an isomorphism for $n \geq l$.
If $\widetilde{M}$ is zero, then $N = 0$.
Thus for any $f \in A_+$ homogeneous with
$\deg(f) + \deg(x) = nd$ and $nd > l$ we see that $fx$ is zero
because $N_{nd} \to N'_{nd}$ and $M_{nd} \to N'_{nd}$ are isomorphisms.
Hence $x$ is irrelevant.
Conversely, assume $M$ is irrelevant. Then $M_{nd}$ is zero
for $n \gg 0$ (see discussion above proposition).
Clearly this implies that $M_{(f_i)} = M_{(f_i^{d/\deg(f_i)})} = 0$,
whence $\widetilde{M} = 0$ by construction.

\medskip\noindent
It follows that the subcategory $\text{Mod}^{fg}_{A, irrelevant}$
is a Serre subcategory of $\text{Mod}^{fg}_A$ as the kernel
of the exact functor $M \mapsto \widetilde M$, see
Homology, Lemma \ref{homology-lemma-kernel-exact-functor}
and Constructions, Lemma \ref{constructions-lemma-proj-sheaves}.
Hence the quotient category on the left of the arrow is defined
in Homology, Lemma \ref{homology-lemma-serre-subcategory-is-kernel}.
To define the functor of the proposition, it suffices to show that
the functor $M \mapsto \widetilde M$ sends irrelevant modules to $0$
which we have shown above.

\medskip\noindent
By Lemma \ref{lemma-coherent-on-proj-general} the proposed quasi-inverse
makes sense. Namely, the lemma shows that
$\mathcal{F} \longmapsto \bigoplus_{n \geq 0} \Gamma(X, \mathcal{F}(n))$
is a functor $\textit{Coh}(\mathcal{O}_X) \to \text{Mod}^{fg}_A$
which we can compose with the quotient functor
$\text{Mod}^{fg}_A \to \text{Mod}^{fg}_A/\text{Mod}^{fg}_{A, irrelevant}$.

\medskip\noindent
By Lemma \ref{lemma-recover-tail-graded-module-general}
the composite left to right to left is isomorphic to the identity functor.
Namely, let $M$ be a finite graded $A$-module and let
$k \in \mathbf{Z}$ sufficiently small and let
$N \to N'$ and $M \to N'$ be as in
Lemma \ref{lemma-recover-tail-graded-module-general}.
Then the kernel and cokernel of $M \to N'$ are nonzero in
only finitely many degrees, hence are irrelevant. Moreover, the
kernel and cokernel of the map $N \to N'$ are zero in all sufficiently
large degrees divisible by $d$, hence these are irrelevant modules too.
Thus $M \to N'$ and $N \to N'$ are both isomorphisms in the quotient
category, as desired.

\medskip\noindent
Finally, let $\mathcal{F}$ be a coherent $\mathcal{O}_X$-module.
Set $M = \bigoplus_{n \in \mathbf{Z}} \Gamma(X, \mathcal{F}(n))$
viewed as a graded $A$-module, so that our functor sends $\mathcal{F}$ to
$M_{\geq 0} = \bigoplus_{n \geq 0} M_n$.
By Properties, Lemma \ref{properties-lemma-proj-quasi-coherent}
the canonical map $\widetilde M \to \mathcal{F}$
is an isomorphism. Since the inclusion map
$M_{\geq 0} \to M$ defines an isomorphism
$\widetilde{M_{\geq 0}} \to \widetilde M$ we conclude that
the composite right to left to right is isomorphic to the identity
functor as well.
\end{proof}

\end{document}
